\chapter{研究原则与语义不变量}

\section{分层维护不同事实}

研究资料包以语义不变量约束写入权限、检索行为和服务表示。下表列出这些不变量及其保护对象。\cite{nous-research-package}

\begin{ReportWideTable}{p{6.4cm}p{9.5cm}}
\captionof{table}{核心语义不变量及其保护对象}\label{tab:invariants}
\toprule
\textbf{不变量} & \textbf{保护原则}\\
\midrule
主体（Subject）$\neq$ 模型（Model） & 主体身份独立于模型名称、接口供应商和会话标识。\\
材料（Artifact）$\neq$ 观察事件（ObservationOccurrence）$\neq$ 记忆（Memory） & 材料内容、主体一次接触和长期认知分别建模。\\
材料提及（Mention）$\neq$ 认知指向（Aboutness） & 分别记录材料中出现的实体与主体认知所指。\\
证据来源谱系 $\neq$ 独立佐证数 & 证据独立性依据来源谱系计算，派生文本保留共同祖先。\\
检索命中 $\neq$ 服务投影 $\neq$ 使用记录 $\neq$ 认知强化 & 检索、呈现和有效使用分别记录；只有合同定义的使用回执可影响可访问性或关联支持度。\\
假设（Hypothesis）$\neq$ 已接受的认知 & 假设经目标领域所有者核验后，方可形成认知修订。\\
目标（Goal）$\neq$ 承诺（Commitment） & 内部期望状态由动机领域维护；对外行为承诺由 Heptalogos 管理。\\
服务投影（Serving Projection）$\neq$ 认知权威（Cognitive Authority） & 中心性、传播势能和向量得分用于检索，不决定长期认知真值。\\
可访问性 $\neq$ 真值 $\neq$ 抑制 $\neq$ 清除 & 回忆资格、主张状态、普通读取限制和管理清除分别维护。\\
自然语言中的“忘记” $\neq$ 删除授权 & 抑制与清除依照类型明确的管理权限执行。\\
\bottomrule
\end{ReportWideTable}

\section{主张类型}

\StatusTag{架构约束}{规定系统应保持的语义关系，例如认知权威与服务投影分离。}

\StatusTag{实验观察}{描述明确试验条件下观察到的性质；局部性质试验与系统级检索质量分别报告。}

\StatusTag{外部依据}{论文或成熟系统帮助界定问题结构，项目数据结构和参数由本项目合同及评测决定。}

\StatusTag{代码事实}{表示指定代码快照中可定位的对象、接口或逻辑；运行行为另以资格记录描述。}

\section{设计取舍约束}

领域所有者由稳定身份、修订、生命周期、写入权限和查询语义界定。成熟库负责通用机制，Nous Wave 负责认知语义。抽象和过程的范围应与承担的责任相称；目标语义明确后，再确定注册、调度和策略机制。

\section{状态标签}

全文使用“必需、目标、已实现、已核验、实验观察、候选、待定”等状态词描述具体能力或主张。代码实现、运行资格和研究结论分别标记；同一算法可已有实现，同时继续作为候选方案参与评测。
